\section{Въведение в уеб-базирани платформи за онлайн търговия на компютърни части}

Настоящата глава поставя разработката в контекста на съществуващите подходи и теоретичните основи на уеб-базираните платформи за електронна търговия. Първо се прави преглед на типичните решения в областта -- от готови SaaS платформи до custom приложения върху общи уеб рамки -- и се обосновава изборът на самостоятелна разработка. След това се въвеждат основните теоретични понятия, върху които стъпва архитектурата на системата: клиент--сървърно взаимодействие, многослойна организация, MVC модел, обектно-релационно моделиране, валидация на данни и сесийно управление.

\subsection{Преглед на съществуващите решения}

За да бъде правилно оценена една софтуерна разработка, тя следва да бъде позиционирана спрямо вече съществуващите решения в областта. Това е особено важно при системите за електронна търговия, тъй като пазарът предлага както готови платформи с богата функционалност, така и множество технологични стекове за изграждане на custom приложения. Настоящият проект не възниква във вакуум; той е реализиран в контекста на вече наложени модели за организация на каталог, поръчка, потребителски профил и административни функции. Целта на настоящия раздел е да очертае този контекст и да покаже с какво разглежданата система се доближава до утвърдените практики и в какво съзнателно се различава.

\subsubsection{Обща картина на съвременните решения за електронна търговия}

Съвременните системи за електронна търговия могат условно да се разделят на три основни групи. Първата група включва напълно управлявани SaaS платформи, при които доставчикът поема голяма част от инфраструктурните и административните задачи. Втората група включва self-hosted системи с готова функционалност и възможност за разширяване чрез плъгини и теми. Третата група обхваща custom решения, разработени върху общопредназначени уеб рамки, при които логиката, структурата на данните и интерфейсът се изграждат специално за конкретния проект.

Изборът между тези подходи зависи от множество фактори: степен на необходима гъвкавост, бюджет, срокове, изисквания към интеграции, организационен контекст и налични технически компетентности. Готовите платформи предлагат бърз старт и богат набор от стандартни функции, но ограничават архитектурния контрол. Custom разработката изисква повече проектантски и програмни усилия, но предоставя възможност системата да бъде моделирана точно според целите на екипа или учебната задача. В този смисъл настоящият проект попада в третата група и следва да бъде разглеждан като контролирана инженерна реализация, а не като конфигуриране на готов продукт.

\subsubsection{Готови платформи: Shopify, WooCommerce и Magento}

Една от най-разпознаваемите SaaS платформи за електронна търговия е Shopify \cite{shopify}. Нейното основно предимство е, че предлага цялостна инфраструктура, готов административен панел, управление на продукти, теми, плащания и интеграции с външни услуги. За бизнеси, които търсят бързо излизане на пазара и не се нуждаят от ниско ниво на контрол над архитектурата, Shopify е разумен избор. От гледна точка на дипломна разработка обаче този тип платформа е по-слабо подходяща, защото голяма част от инженерните решения са вече предварително взети от доставчика, а възможностите за задълбочен архитектурен анализ са ограничени.

WooCommerce \cite{woocommerce} представлява различен подход. То е разширение за WordPress, което превръща система за управление на съдържание в електронен магазин. Предимството му е в лесната инсталация, голямата екосистема от теми и плъгини и ниската бариера за въвеждане. Подходът е особено полезен, когато вече съществува WordPress сайт и електронният магазин трябва да бъде интегриран към него. Недостатъкът е, че архитектурата на приложението остава силно зависима от структурата на WordPress и от начина, по който плъгините взаимодействат помежду си. Това ограничава чистотата на модела и прозрачността на вътрешната бизнес логика.

Magento Open Source \cite{magento} е пример за мощна self-hosted платформа, ориентирана към по-сложни сценарии, по-големи каталози и по-високи изисквания към конфигурация и разширяване. Тя предлага богат набор от модули, административни възможности и интеграционни точки. В същото време сложността ѝ е значително по-висока от тази на WooCommerce или SaaS решенията. За малка или учебна разработка такава платформа често е прекалено тежка, тъй като значителна част от усилието се пренасочва към усвояване на рамката и нейните вътрешни механизми, вместо към ясно проектно моделиране на конкретното приложение.

Съпоставката между тези платформи показва, че готовите решения са силни, когато бизнесът има нужда от бърза operational готовност. Те обаче не са оптимални, когато основната цел е демонстрация и анализ на фундаментални принципи от софтуерното инженерство. Именно поради това в настоящия проект е предпочетена самостоятелна разработка върху обща уеб рамка.

\begin{table}[H]
    \centering
    \caption{Съпоставка между типични подходи за изграждане на електронен магазин}
    \label{tab:ecommerce-platforms}
    \begin{tabularx}{\textwidth}{L{2.7cm}YYY}
        \toprule
        Подход & Предимства & Ограничения & Подходящост за дипломна разработка \\
        \midrule
        Shopify & Бърз старт, хоствана инфраструктура, готови интеграции & Ограничен архитектурен контрол, зависимост от доставчик & Ниска, защото инженерните решения са силно абстрахирани \\
        WooCommerce & Лесно внедряване, богата екосистема, ниска бариера за вход & Силна зависимост от WordPress, натрупване на плъгини & Средна, ако фокусът е върху конфигурация, не върху архитектура \\
        Magento & Богата функционалност, разширяемост, зрял търговски модел & Висока сложност и тежка среда за внедряване & Средна към ниска при ограничен времеви ресурс \\
        Custom Spring приложение & Пълен контрол върху модел, код и архитектура & По-голяма първоначална разработка & Висока, защото позволява анализ на всички слоеве \\
        \bottomrule
    \end{tabularx}
\end{table}

Таблица \ref{tab:ecommerce-platforms} показва, че custom Spring приложение не е най-бързият път към готова търговска платформа, но е най-подходящият избор, когато целта е демонстрация на архитектурна последователност, моделиране на предметна област и проследимост на решенията от потребителския интерфейс до базата данни.

\subsubsection{Custom уеб приложения и ролята на Spring екосистемата}

В Java екосистемата Spring Boot е сред най-често използваните средства за бързо изграждане на самостоятелни уеб приложения \cite{springboot}. Неговото основно предимство се състои в стандартизирания модел за конфигурация, автоматична инициализация на инфраструктурата и лесна интеграция между различни модули като MVC, сигурност, достъп до данни, шаблони и външни услуги. Това го прави подходяща основа за приложения, които трябва да имат ясен архитектурен гръбнак, но да не се обременяват от прекомерна ръчна конфигурация.

Spring MVC \cite{springmvc} е особено важен в контекста на настоящата разработка, защото предлага естествен модел за маршрутизация на HTTP заявки, binding на параметри, интеграция с валидация и лесна връзка с шаблонни енджини като Thymeleaf \cite{thymeleaf}. По този начин всеки URL маршрут може ясно да бъде свързан с контролер, бизнес операция и визуализация. Това е съществено предимство при дипломна разработка, защото позволява по-ясно описание на потоците в системата.

Spring Data JPA \cite{springdatajpa} също играе важна роля. Вместо значителна част от времето да бъде изразходвана за ръчно писане на SQL слой и трансформации между резултатни множества и обекти, JPA моделът позволява домейн обектите да бъдат описани като entity класове с релации, а repository интерфейсите да делегират част от рутинната инфраструктурна логика на рамката. Това не премахва нуждата от добро моделиране, но ускорява реализацията на устойчив слой за съхранение на данни.

За малки и средни приложения изборът на Spring Boot носи и друго предимство: позволява решението да остане монолитно в архитектурен смисъл, но вътрешно добре разделено по модули. Това е особено подходящо за учебен проект, при който е по-важно вътрешната организация да бъде прозрачна и последователна, отколкото приложението да бъде преждевременно разделяно в множество услуги.

\subsubsection{Архитектурни подходи при уеб приложенията}

Освен между различни платформи, изборът трябва да бъде направен и между различни архитектурни парадигми на самото уеб приложение. В най-общ план могат да бъдат разграничени два доминиращи модела. Първият е сървърно рендираното приложение, при което HTML се генерира на сървъра и се изпраща към клиента като готова страница. Вторият е клиентски ориентираният модел, при който front-end приложението е самостоятелно, а сървърът предоставя основно REST или GraphQL API.

Сървърно рендираният модел има няколко съществени предимства за разглеждания проект. Той намалява общата архитектурна сложност, тъй като представящият слой и бизнес логиката остават в една и съща технологична среда. Освен това позволява бърза реализация на форми, валидации и навигационни потоци, без да е необходимо изграждането на отделно front-end приложение с независим процес на компилация и разгръщане. За дипломен проект, чиято основна цел е да покаже цялостен back-end и UI процес с ограничен технологичен риск, това е рационален избор.

Клиентски ориентираният модел, типичен за приложения, базирани на React, Angular или Vue, предоставя по-голяма динамичност на интерфейса и по-естествена интеграция с мобилни и външни клиенти. Той обаче изисква отделна front-end архитектура, по-ясно API контрактуване, повече внимание към клиентското състояние и често по-сложен модел за сигурност и автентикация. В контекста на текущата задача това би довело до разширяване на технологичния обхват за сметка на фокуса върху систематичната реализация на основните бизнес процеси.

Настоящият проект избира хибриден, но доминиращо сървърно рендиран подход. Основната функционалност се визуализира чрез Thymeleaf, докато отделна REST крайна точка е използвана за демонстрация на програмно достъпна бизнес информация. По този начин системата запазва простота, но същевременно показва възможност за комбиниране на HTML и JSON интерфейси.

\subsubsection{Позициониране на настоящия проект}

Разработваната система следва да бъде разбрана като custom монолитно приложение, ориентирано към учебно-приложна стойност. То не се конкурира пряко с пълнофункционални търговски платформи, а по-скоро демонстрира как чрез добре подбрани технологии могат да бъдат реализирани основните механизми на електронен магазин в прозрачна и разбираема архитектурна форма.

Проектът заема междинна позиция между два екстрема. От едната страна са платформите тип „всичко готово“, които скриват значителна част от инженерния процес. От другата страна са големите enterprise решения с микросервизна или силно модулна архитектура, чиято сложност е непропорционална на целите на дипломната задача. Настоящата система избира балансиран път: достатъчно реалистична функционалност, реална база данни, реални HTML шаблони, реална e-mail интеграция и реална структура на сорс кода, но без прекомерни инфраструктурни усложнения.

Тази позиция я прави особено подходяща за академичен анализ. В нея могат ясно да бъдат проследени домейн обектите, заявките към системата, маршрутизацията, ролята на валидацията, управлението на сесията, трансформациите между модели и мястото на конфигурацията. Именно тази проследимост е едно от основните предимства на избрания подход.

\subsubsection{Изводи от обзорната част}

Прегледът на съществуващите решения показва, че в областта на електронната търговия няма универсално най-добър подход; подходящият избор зависи от целите на конкретния проект. Ако основната задача е бързо въвеждане в експлоатация, SaaS и плъгин-базираните решения имат ясни предимства. Ако обаче водеща е необходимостта от архитектурна прозрачност, контрол върху бизнес логиката и възможност за детайлно инженерно описание, custom уеб приложение върху Spring екосистемата е по-подходящ избор.

Именно в този смисъл настоящата дипломна разработка е обоснована: тя не се стреми да възпроизведе в пълнота зрелите търговски платформи, а да изгради представителен, работещ и аналитично обясним модел на онлайн магазин.

\subsection{Теоретични основи}

За да бъде разбрана и аргументирана архитектурата на разглежданата система, е необходимо да бъдат въведени няколко основни теоретични понятия. Макар проектът да има подчертано практическа насоченост, неговата реализация стъпва върху утвърдени концепции от областта на уеб инженерството и софтуерната архитектура: клиент--сървърно взаимодействие, многослойна организация, MVC модел, обектно-релационно моделиране, управление на сесия и валидация на данни. Този раздел не цели изчерпателен теоретичен обзор, а по-скоро въвежда именно онези идеи, които са пряко необходими за разбиране на последващото описание на системата.

\subsubsection{Клиент--сървърна природа на уеб приложението}

В основата на всяка уеб система стои клиент--сървърният модел. Клиентът, най-често уеб браузър, инициира HTTP заявка към сървъра. Сървърът интерпретира заявката, прилага бизнес логика, извлича или записва данни и връща отговор, който може да бъде HTML документ, пренасочване, файл или структуриран формат като JSON. Тази обща схема е в сила и за разглежданата система: браузърът на потребителя изпраща заявки към контролерите на приложението, а те в зависимост от маршрута връщат страници, извършват пренасочване или генерират REST отговор.

От инженерна гледна точка важна е идеята, че протоколът HTTP е по своята природа безсъстояниен. Това означава, че отделните заявки са независими една от друга, освен ако приложението не въведе собствен механизъм за проследяване на потребителския контекст. В настоящия проект тази роля се поема от \texttt{HttpSession}, където се съхранява електронната поща на текущо удостоверения потребител. Така системата постига достатъчно прост модел на сесийно състояние, подходящ за учебно приложение, без да въвежда по-сложни механизми като токени, филтри и централизирана сигурност.

\subsubsection{Многослойна архитектура}

Многослойната архитектура е един от най-устойчивите организационни модели в enterprise разработката \cite{fowler,sommerville}. Нейната идея е различните отговорности в системата да бъдат разделени на слоеве с относително ясни граници. В типичния случай могат да бъдат разграничени визуален слой, приложен слой за управление на заявки и потоци, бизнес слой за домейн логика и слой за достъп до данни.

Предимството на този модел е, че промени в един слой могат да бъдат извършвани с по-малко странични ефекти върху останалите. Например промяна в HTML шаблон не би следвало да налага промяна в entity класовете, а промяна в начина на съхранение на данни не би следвало да засяга непосредствено логиката на формите. Разбира се, в реални системи винаги има определено взаимодействие между слоевете, но многослойният модел създава дисциплина, която намалява вероятността от силно преплитане на отговорности.

В настоящия проект многослойността е реализирана чрез контролери, услуги, repository интерфейси и шаблони. Този подход е особено подходящ за образователна разработка, защото позволява лесно да се проследи движението на данните от потребителската форма към persistence слоя и обратно.

\subsubsection{Моделът MVC}

Model--View--Controller е класически архитектурен шаблон за уеб приложения и графични системи \cite{springmvc}. В контекста на Spring MVC той се проявява по следния начин:

\begin{itemize}[leftmargin=*]
    \item \textbf{Model} представя данните, които се подават към визуализацията;
    \item \textbf{View} е шаблонът, който визуализира тези данни пред потребителя;
    \item \textbf{Controller} обработва входящата заявка и избира коя view компонента да бъде върната.
\end{itemize}

Важно е да се подчертае, че в Spring MVC терминът „модел“ не съвпада напълно с домейн модела на приложението. В конкретната реализация домейн обектите са entity класовете, а обектът \texttt{Model} в контролерите представлява контейнер за атрибути, които ще бъдат достъпни от Thymeleaf шаблона. Това разграничение е съществено, защото показва, че MVC не е заместител на домейн модела, а механизъм за организиране на взаимодействието между потребителската заявка и визуализирания резултат.

В разглежданата система MVC моделът е видим почти във всеки маршрут. Контролерът получава заявка за регистрация, логин, каталог, количка или профил; извиква съответна услуга; поставя резултата в модела; и връща име на шаблон. Тази последователност е една от причините проектът да бъде добър пример за ясна учебна архитектура.

\begin{figure}[H]
    \centering
    \begin{tikzpicture}
        \node[appblock] (browser) {Уеб браузър};
        \node[appblock, below=1.4cm of browser] (controller) {Controller};
        \node[appblock, below left=1.5cm and 2.4cm of controller] (service) {Service слой};
        \node[appblock, below right=1.5cm and 2.4cm of controller] (view) {Thymeleaf view};
        \node[appblock, below=1.4cm of service] (repo) {Repository слой};
        \node[appblock, below=1.4cm of repo] (db) {MySQL база данни};

        \draw[appline] (browser) -- node[right]{HTTP заявка} (controller);
        \draw[appline] (controller) -- node[left]{бизнес логика} (service);
        \draw[appline] (service) -- (repo);
        \draw[appline] (repo) -- (db);
        \draw[appline] (controller) -- node[right]{модел} (view);
        \draw[appline] (view) |- node[pos=0.25,right]{HTML} (browser);
    \end{tikzpicture}
    \caption{Обобщен поток на обработка при сървърно рендирана заявка}
    \label{fig:mvc-request-flow}
\end{figure}

Фигура \ref{fig:mvc-request-flow} илюстрира основния път на заявката в системата. Макар той да изглежда линеен, зад всяка стъпка стои отделен набор от отговорности: контролерите управляват маршрутизацията, услугите формулират правилата на домейна, repository компонентите делегират работата с базата данни, а шаблоните осигуряват потребителската визуализация.

\subsubsection{Dependency Injection и конфигурация на компоненти}

Една от основните идеи в Spring е \textit{dependency injection}, тоест доставяне на необходимите зависимости от външен контейнер, вместо създаването им на ръка във всеки клас. Практически това позволява класовете да бъдат по-слабо свързани, по-лесно тествани и по-ясно конфигурирани \cite{springboot}. Когато дадена услуга зависи от repository интерфейс, model mapper или password encoder, тя ги получава чрез конструктора си, вместо сама да ги инстанцира.

В текущия проект това е реализирано последователно: контролерите приемат необходимите услуги и инфраструктурни компоненти чрез конструкторна инжекция, а конфигурационният клас дефинира bean инстанции за \texttt{ModelMapper} и \texttt{PasswordEncoder}. Това решение е важно не само за чистотата на кода, но и за архитектурната проследимост: когато се разглежда даден клас, веднага се вижда от кои други компоненти зависи и каква роля изпълнява.

\subsubsection{Домейн модел и обектно-релационно съответствие}

Информационните системи често трябва да представят едни и същи понятия в два различни свята: обектния свят на програмния език и релационния свят на базата данни. Обектно-релационното съответствие се стреми да преодолее тази разлика чрез механизъм, при който класовете и връзките между тях се свързват с таблици, колони и външни ключове \cite{springdatajpa}. В Java екосистемата това обикновено се реализира чрез JPA анотации върху entity класове.

В разглежданата система класове като \texttt{User}, \texttt{Company}, \texttt{Product}, \texttt{ShoppingCartEntity} и \texttt{HistoryEntity} са маркирани като entity обекти и описват основните релации в домейна. Така се постига сравнително естествена връзка между предметната област и persistence слоя. Например един потребител има множество фирми, количка и история на поръчки; продуктът принадлежи към фирма; а записът в количката свързва конкретен продукт с конкретен потребител и желано количество.

Тук е важно да се направи разграничение между \textit{домейн модел} и \textit{модели за пренос на данни}. В проекта не всички данни се обменят чрез entity обектите. За някои операции са въведени binding модели за входни форми, service модели за вътрешен трансфер между слоеве и view модели за изход към интерфейса. Това е добър архитектурен избор, тъй като ограничава директното изтичане на persistence структурата към потребителския интерфейс.

\subsubsection{Валидация на входни данни}

Въвеждането на данни от потребител е една от най-честите причини за логически грешки, нарушаване на ограничения в базата или несъответствие между очакваните и действителните стойности. Поради това валидацията на входните данни е фундаментална характеристика на всяка надеждна информационна система. В Spring приложенията това често се реализира чрез Java Bean Validation анотации и обработка на \texttt{BindingResult} \cite{hibernatevalidator}.

В конкретния проект ограниченията са дефинирани основно в binding моделите. За полета като име, адрес, количество, цена и други са зададени правила от типа минимална и максимална дължина, непразна стойност или положително число. Този подход има няколко предимства. Първо, позволява формалното описание на допустимите входни стойности да бъде концентрирано на едно място. Второ, дава възможност контролерът да реагира систематично при невалиден вход, като върне потребителя към съответната форма с индикация за грешка. Трето, намалява риска в service слоя да достигат невалидни обекти.

Въпреки това е важно да се отбележи, че валидацията в текущия проект е частична. Например някои binding модели не съдържат богато множество от ограничения, а част от логическите проверки се извършват на по-късен етап в контролерите или услугите. Това е приемливо за учебен проект, но при по-зряла система би било полезно валидационните правила да бъдат още по-последователни и формализирани.

\subsubsection{Управление на потребителска идентичност и защита на паролите}

Системите за електронна търговия оперират с чувствителни данни и поради това дори базовият модел на автентикация трябва да бъде проектиран разумно. Един от най-важните принципи е паролите никога да не се съхраняват в явен вид, а да бъдат обработвани чрез устойчив механизъм за хеширане \cite{springsecuritycrypto,owasp}. В текущия проект тази роля се изпълнява от \texttt{Pbkdf2PasswordEncoder}, който осигурява по-добра защита от директно съхраняване или използване на бързи, неподходящи за пароли хеш функции.

Моделът на удостоверяване в проекта е сравнително опростен. При успешен вход електронната поща на потребителя се записва в \texttt{HttpSession}. От тази стойност впоследствие се определя дали дадена заявка е позволена и кой е текущият потребител. Този подход е лесен за разбиране и достатъчен за малка учебна система, но има ограничена изразителност спрямо по-пълноценни решения със сигурностен контекст, роли, филтри и декларативни правила за достъп. По-късно в документа това ще бъде разгледано като естествена възможност за бъдещо надграждане.

\subsubsection{Сървърно рендирани шаблони и локализация}

Thymeleaf е шаблонен механизъм за сървърно рендиране, който интегрира естествено HTML маркировка и изрази за работа с данни от модела \cite{thymeleaf}. Неговото предимство е, че прави HTML шаблоните сравнително четими и при липса на изпълнение, а в същото време позволява условна визуализация, итериране по колекции, двупосочно свързване на форми и използване на локализирани съобщения.

В настоящия проект Thymeleaf се използва като основен механизъм за изграждане на потребителския интерфейс. Формите за регистрация, вход, добавяне на компания, добавяне на продукт, профил, каталог, количка и история са реализирани именно по този начин. Допълнителен теоретичен аспект тук е локализацията. Чрез message bundles и интерцептор за смяна на езика системата демонстрира базов модел за поддръжка на повече от един език, като конкретно предлага английски и български вариант на част от интерфейсните текстове.

\subsubsection{REST интерфейс и представяне на данни}

Макар приложението да е доминиращо сървърно рендирано, в него е наличен и примерен REST интерфейс, връщащ информация за максималната цена на продукт. Теоретичното значение на тази част е, че показва как една и съща бизнес логика може да бъде използвана през различни форми на представяне. Вместо HTML, при REST сценария резултатът се сериализира до JSON и може да бъде използван от JavaScript код в браузъра или от външен клиент. Подобен стил на взаимодействие е в съзвучие с идеите на уеб архитектурата и ресурсно ориентираното представяне на данни \cite{fielding}.

Дори този относително малък REST пример има методологична стойност. Той показва, че границата между традиционно MVC приложение и API ориентирано приложение не е абсолютно строга. Една добре структурирана система може да комбинира двата подхода и да предоставя различни изгледи върху една и съща домейн логика според нуждите на клиента.

\subsubsection{Обобщение на теоретичната рамка}

Разгледаните понятия очертават теоретичния фундамент на настоящата разработка. Клиент--сървърният модел определя начина, по който потребителят взаимодейства със системата. Многослойната архитектура организира вътрешните ѝ отговорности. MVC моделът структурира уеб заявките и визуализациите. Обектно-релационното моделиране свързва домейн обектите с базата данни. Валидацията ограничава невалидния вход. Хеширането на пароли и управлението на сесия осигуряват базова сигурност. Thymeleaf и локализацията оформят представящия слой, а REST крайна точка показва алтернативен начин за експониране на данни.

Тази теоретична рамка не е само фон; тя е пряко материализирана в структурата на проекта. Следващата глава преминава от общите понятия към конкретното проектиране на уеб-базираната платформа -- изискванията, архитектурата и модела на данните.
